\section{Выводы}

В ходе лабораторной работы были экспериментально подтверждены ключевые аспекты поляризации света.  

\begin{itemize}
    \item \textbf{Угол Брюстера} позволил определить коэффициент преломления эбонита (\( n \approx 1.43–1.54 \)), близкий к табличному значению (\( 1.6 \)). Погрешности связаны с неточностями установки угла, что подчеркивает важность стабильности оптической системы.  

    \item \textbf{Поляризация отражённого и преломлённого света}:
    \begin{itemize}
        \item Отражённый свет при угле Брюстера полностью поляризован в плоскости падения.  
        \item Преломлённый свет частично поляризован перпендикулярно плоскости, причём стопа пластинок усиливает этот эффект.  
    \end{itemize}

    \item \textbf{Работа с пластинками}:
    \begin{itemize}
        \item Пластинка \( \lambda/4 \) преобразует линейную поляризацию в круговую, а \( \lambda/2 \) поворачивает плоскость колебаний на \( 180^\circ \).  
        \item Определение «быстрой» и «медленной» осей подтвердило зависимость цвета от разности фаз и ориентации пластинок.  
    \end{itemize}

    \item \textbf{Интерференция} в мозаичной слюдяной пластинке продемонстрировала, как разная толщина элементов создаёт цветовые эффекты. Вращение анализатора меняло проекцию вектора \( \mathbf{E} \), влияя на интерференционную картину.  

    \item \textbf{Направление вращения вектора \( \mathbf{E} \)} в эллиптически поляризованном свете определялось через сдвиг фаз. При \( \Delta\phi = \pi/2 \) вращение происходило против часовой стрелки, что согласуется с теорией.  
\end{itemize}

Работа показала, как экспериментальные методы (использование поляроидов, пластинок \( \lambda/4 \), \( \lambda/2 \), анализ интерференции) позволяют исследовать свойства поляризованного света и подтверждать теоретические модели.